\chapter{จุลินทรีย์กับการควบคุมศัตรูพืชโดยชีววิธี}
\label{chap:ch04}

% ==========================================================
% แผนบริหารการสอนประจำบทที่ 4
% ==========================================================
\section*{แผนบริหารการสอนประจำบทที่ 4}
\addcontentsline{toc}{section}{แผนบริหารการสอนประจำบทที่ 4}

\noindent\textbf{หัวข้อเนื้อหาประจำบท}
\begin{enumerate}
    \item นิยามและหลักการของการควบคุมศัตรูพืชโดยชีววิธี (Biological Control)
    \item กลไกหลักของจุลินทรีย์ปฏิปักษ์ (การเป็นปรสิต การหลั่งสารต้านจุลชีพ และการชักนำภูมิคุ้มกัน)
    \item ชีววิทยาและการประยุกต์ใช้เชื้อราไตรโคเดอร์มา (\textit{Trichoderma} spp.)
    \item แบคทีเรียปฏิปักษ์กลุ่มบาซิลลัส (\textit{Bacillus} spp.) และแอกทิโนมัยซีท
    \item จุลินทรีย์ก่อโรคแมลงศัตรูพืช (แบคทีเรียบีที เชื้อราบิวเวอเรีย และไวรัสเอ็นพีวี)
\end{enumerate}

\noindent\textbf{วัตถุประสงค์เชิงพฤติกรรม}
\begin{enumerate}
    \item อธิบายกลไกทางสรีรวิทยาและชีวเคมีของจุลินทรีย์ปฏิปักษ์ในการยับยั้งเชื้อโรคพืชได้
    \item ระบุชนิดของจุลินทรีย์ควบคุมชีววิธีและจับคู่กับโรคหรือแมลงศัตรูพืชเป้าหมายได้อย่างแม่นยำ
    \item อธิบายกลไกการออกฤทธิ์ของสารพิษ Cry protein จากแบคทีเรีย \textit{Bacillus thuringiensis} ได้
    \item ผลิตและขยายหัวเชื้อจุลินทรีย์ปฏิปักษ์บนเมล็ดธัญพืชตามหลักวิชาการได้
\end{enumerate}

\noindent\textbf{กิจกรรมการเรียนการสอน}
\begin{enumerate}
    \item การบรรยายและวิเคราะห์กลไก Dual Culture Assay ระหว่างจุลินทรีย์ปฏิปักษ์กับเชื้อราก่อโรค
    \item ปฏิบัติการเพาะเลี้ยงเชื้อราไตรโคเดอร์มาบนเมล็ดข้าวสุกและตรวจสอบความบริสุทธิ์
    \item การทดสอบประสิทธิภาพของสารสกัดชีวภาพในการยับยั้งเส้นใยเชื้อราสาเหตุโรครากเน่า
\end{enumerate}

\noindent\textbf{สื่อการเรียนการสอน}
\begin{enumerate}
    \item เอกสารคำสอนบทที่ 4 และภาพประกอบแผนผังกลไกชีววิธี
    \item ตัวอย่างจานอาหารเลี้ยงเชื้อ Dual Culture แสดงวงยับยั้ง (Inhibition Zone)
\end{enumerate}

\noindent\textbf{การประเมินผล}
\begin{enumerate}
    \item การตรวจประเมินคุณภาพหัวเชื้อสดไตรโคเดอร์มาที่นักศึกษาผลิตในชั่วโมงปฏิบัติการ
    \item การทดสอบความรู้ความเข้าใจเกี่ยวกับกลไก ISR และการออกฤทธิ์ทางชีวภาพ
\end{enumerate}

\newpage

% ==========================================================
% เนื้อหาบทเรียน
% ==========================================================

\section{หลักการควบคุมศัตรูพืชโดยชีววิธี}
\index{การควบคุมโดยชีววิธี (Biological Control)}

\textbf{การควบคุมศัตรูพืชโดยชีววิธี (Biological Control)} หมายถึง การใช้สิ่งมีชีวิตหรือผลผลิตจากสิ่งมีชีวิตในการควบคุมและลดประชากรของสิ่งมีชีวิตก่อโรคหรือแมลงศัตรูพืชให้อยู่ในระดับที่ไม่ก่อให้เกิดความเสียหายทางเศรษฐกิจ

การควบคุมโดยชีววิธีมีข้อได้เปรียบเหนือการใช้สารเคมีสังเคราะห์หลายประการ ได้แก่ มีความจำเพาะเจาะจงต่อสิ่งมีชีวิตเป้าหมายสูง ไม่ตกค้างสะสมในผลผลิตและสิ่งแวดล้อม ไม่ทำลายแมลงผสมเกสรและสิ่งมีชีวิตที่เป็นประโยชน์ และไม่กระตุ้นให้ศัตรูพืชเกิดการกลายพันธุ์สร้างความต้านทานต่อยาอย่างรวดเร็ว

\section{กลไกของจุลินทรีย์ปฏิปักษ์ในการควบคุมโรคพืช}
\index{จุลินทรีย์ปฏิปักษ์ (Antagonistic Microorganisms)}

จุลินทรีย์ปฏิปักษ์ควบคุมโรคพืชผ่าน 3 กลไกหลักที่ทำงานประสานกัน (Tripartite Mechanism) ได้แก่

\subsection{1. การเป็นปรสิตบนเส้นใยรา (Mycoparasitism)}
\index{การเป็นปรสิตบนเส้นใยรา (Mycoparasitism)}
เป็นกระบวนการที่เชื้อราปฏิปักษ์ เช่น \textit{Trichoderma} ตรวจจับสารสัญญาณเคมีจากเชื้อราก่อโรค (Chemotropism) เจริญเส้นใยเข้าไปแนบชิดและพันรัดรอบเส้นใยของเชื้อราเป้าหมาย จากนั้นจะสร้างโครงสร้างคล้ายปุ่มดูด (Appressorium-like structures) และหลั่งเอนไซม์ย่อยสลายผนังเซลล์ (Cell Wall Degrading Enzymes หรือ CWDEs) ได้แก่ เอนไซม์ไคติเนส (Chitinase) เบตา-1,3-กลูคาเนส ($\beta$-1,3-glucanase) และโพรทีเอส (Protease) เพื่อเจาะทะลุผนังเซลล์และดูดกินสารอาหารภายในจนเซลล์เชื้อก่อโรคแตกสลาย

\subsection{2. การสร้างสารปฏิชีวนะและการแก่งแย่ง (Antibiosis and Competition)}
\index{การสร้างสารปฏิชีวนะ (Antibiosis)}
จุลินทรีย์ปฏิปักษ์หลั่งสารเมแทบอไลต์ทุติยภูมิ (Secondary Metabolites) ที่มีฤทธิ์ยับยั้งหรือฆ่าเชื้อราและแบคทีเรียก่อโรค เช่น เชื้อราไตรโคเดอร์มาผลิตสารกลุ่มเพปไทโบโลส์ (Peptaibols), วิริดิน (Viridin) และ 6-pentyl-$\alpha$-pyrone (6-PP) ซึ่งมีกลิ่นหอมคล้ายมะพร้าว ส่วนแบคทีเรีย \textit{Bacillus subtilis} ผลิตสารลิโปเปปไทด์ (Lipopeptides) เช่น เซอร์แฟกติน (Surfactin), อิทูริน (Iturin) และเฟนจิซิน (Fengycin) ซึ่งแทรกตัวเข้าสู่เยื่อหุ้มเซลล์ของเชื้อราก่อโรคทำให้เกิดรูรั่วและตายในที่สุด นอกจากนี้ การเจริญเติบโตอย่างรวดเร็วและการสร้างไซเดอโรฟอร์เพื่อแย่งชิงธาตุเหล็กยังช่วยตัดวงจรอาหารของเชื้อก่อโรคอีกด้วย

\subsection{3. การชักนำให้พืชสร้างความต้านทานโรค (Induced Systemic Resistance หรือ ISR)}
\index{การชักนำความต้านทานโรค (ISR)}
การเข้าครอบครองพื้นที่รากของจุลินทรีย์ปฏิปักษ์จะส่งสัญญาณเตือนภัยทางชีวภาพกระตุ้นระบบภูมิคุ้มกันของพืชผ่านวิถีกรดแจสโมนิก (Jasmonic acid) และเอทิลีน (Ethylene) ทำให้พืชสร้างสารต้านทานโรคทั่วทั้งต้น เช่น สารไฟโตอเล็กซิน (Phytoalexins) และสะสมลิกนินทำให้ผนังเซลล์พืชหนาขึ้น

\begin{figure}[H]
\centering
\begin{tikzpicture}[scale=0.95]
    % Central circular node
    \node (core) [circle, draw=agriGreen, fill=agriMint, line width=2pt, minimum size=3.6cm, align=center, font=\small\bfseries\color{agriGreen}] 
        {กลไกหลัก\\ชีววิธี\\(Biocontrol\\Mechanisms)};

    % Top node: Mycoparasitism
    \node (p1) [rectangle, rounded corners=8pt, draw=agriLeaf, fill=agriSoftBg, line width=1.4pt, minimum width=4.5cm, minimum height=1.4cm, align=center, font=\footnotesize, above=1.8cm of core]
        {\textbf{1. การเป็นปรสิต (Mycoparasitism)}\\พันรัดเส้นใย หลั่งเอนไซม์ Chitinase,\\$\beta$-1,3-glucanase ย่อยผนังเซลล์};

    % Bottom Left node: Antibiosis
    \node (p2) [rectangle, rounded corners=8pt, draw=agriGold, fill=yellow!15, line width=1.4pt, minimum width=4.5cm, minimum height=1.4cm, align=center, font=\footnotesize, below left=1.6cm and 0.2cm of core]
        {\textbf{2. สารปฏิชีวนะ (Antibiosis)}\\หลั่ง Lipopeptides, Peptaibols\\แย่งจับเหล็กด้วย Siderophore};

    % Bottom Right node: ISR
    \node (p3) [rectangle, rounded corners=8pt, draw=agriBlue, fill=agriBlue!10, line width=1.4pt, minimum width=4.5cm, minimum height=1.4cm, align=center, font=\footnotesize, below right=1.6cm and 0.2cm of core]
        {\textbf{3. ชักนำภูมิคุ้มกัน (ISR)}\\กระตุ้น Jasmonic Acid \& Ethylene\\สะสม Phytoalexins และ Lignin};

    % Connecting lines with arrows
    \draw[-{Stealth[scale=1.3]}, line width=1.6pt, color=agriLeaf] (core) -- (p1);
    \draw[-{Stealth[scale=1.3]}, line width=1.6pt, color=agriGold] (core) -- (p2);
    \draw[-{Stealth[scale=1.3]}, line width=1.6pt, color=agriBlue] (core) -- (p3);
\end{tikzpicture}
\caption{แผนผังกลไก 3 ทิศทางของการควบคุมเชื้อโรคพืชโดยชีววิธี (Tripartite Biocontrol Mechanisms)}
\label{fig:biocontrol_mechanisms}
\end{figure}

\section{จุลินทรีย์ปฏิปักษ์พื้นถิ่นเพื่อการควบคุมโรคไม้ผลในภาคตะวันออก}
\index{แอกติโนมัยซีทปฏิปักษ์ (Antagonistic Actinomycetes)}
\index{การควบคุมโรครากเน่าโคนเน่า (Root Rot Biocontrol)}

การระบาดของโรคพืชสำคัญในสวนผลไม้ภาคตะวันออก เช่น โรครากเน่าโคนเน่าในทุเรียนที่เกิดจากเชื้อ \textit{Phytophthora palmivora} และโรคแอนแทรคโนสในผลไม้เมืองร้อนที่เกิดจากเชื้อ \textit{Colletotrichum} spp. ก่อให้เกิดความเสียหายทางเศรษฐกิจอย่างมหาศาล การวิจัยทางจุลชีววิทยาการเกษตรของมหาวิทยาลัยราชภัฏรำไพพรรณีจึงมุ่งเน้นการค้นหาจุลินทรีย์ปฏิปักษ์สายพันธุ์ท้องถิ่นที่มีความสามารถในการยับยั้งเชื้อก่อโรคได้อย่างมีเสถียรภาพ ดังนี้

\subsection{การคัดกรองแบคทีเรียปฏิปักษ์ต้านเชื้อราจากดินสวนผลไม้จันทบุรี}
Chanthamalee and Puckdee (2023) ได้ดำเนินการสำรวจและคัดแยกแบคทีเรียปฏิปักษ์จากดินสวนผลไม้ในจังหวัดจันทบุรี และทดสอบประสิทธิภาพการยับยั้งเชื้อราก่อโรคพืชด้วยวิธีเลี้ยงเชื้อร่วม (Dual Culture Assay) ผลการทดสอบพบแบคทีเรียสายพันธุ์ท้องถิ่นที่มีศักยภาพสูงในการยับยั้งการเจริญของเส้นใยราก่อโรคได้อย่างเด่นชัด โดยสร้างวงยับยั้งการเจริญ (Inhibition Zone) ที่กว้าง และผลิตสารเมแทบอไลต์นอกเซลล์ที่ทนความร้อน ซึ่งช่วยลดการงอกของสปอร์เชื้อราได้อย่างมีนัยสำคัญ

\subsection{แอกติโนมัยซีทปฏิปักษ์ \textit{Streptomyces} sp. SR1301 จากพื้นที่ปกปักทรัพยากร อพ.สธ.}
Puckdee and Chanthamalee (2026) ประสบความสำเร็จในการแยกและทดสอบเบื้องต้นถึงฤทธิ์ยับยั้งแบคทีเรียของเชื้อ \textit{Streptomyces} sp. SR1301 ที่แยกได้จากตัวอย่างดินในพื้นที่โครงการอนุรักษ์พันธุกรรมพืชอันเนื่องมาจากพระราชดำริฯ (อพ.สธ.) มหาวิทยาลัยราชภัฏรำไพพรรณี

เชื้อ \textit{Streptomyces} sp. SR1301 สร้างสารออกฤทธิ์ทางชีวภาพที่มีขอบเขตการต้านจุลชีพกว้างขวาง สามารถยับยั้งการเจริญเติบโตของทั้งแบคทีเรียและเชื้อราก่อโรคพืช โดยสารสกัดหยาบจากน้ำเลี้ยงเชื้อแสดงประสิทธิภาพการทำลายเยื่อหุ้มเซลล์ของเชื้อก่อโรค การค้นพบสายพันธุ์ SR1301 จากพื้นที่ปกปักพันธุกรรมสะท้อนถึงคุณค่าของความหลากหลายทางชีวภาพจุลินทรีย์พื้นถิ่นในการนำมาพัฒนาเป็นชีวภัณฑ์ป้องกันกำจัดศัตรูพืชที่ปลอดภัยและยั่งยืน

\section{จุลินทรีย์ควบคุมแมลงศัตรูพืช}
\index{จุลินทรีย์ควบคุมแมลง (Entomopathogenic Microorganisms)}

การควบคุมแมลงศัตรูพืชด้วยจุลินทรีย์อาศัยเชื้อโรคเฉพาะเจาะจงที่ก่อให้เกิดโรคตายในแมลง ได้แก่
\begin{enumerate}
    \item \textbf{แบคทีเรียบีที (\textit{Bacillus thuringiensis} หรือ Bt)} ในระหว่างกระบวนการสร้างสปอร์ แบคทีเรียบีทีจะสร้างผลึกสารพิษโปรตีนรูปเพชร เรียกว่า \textbf{$\delta$-endotoxin} หรือ \textbf{Cry protein} เมื่อแมลงศัตรูพืช เช่น หนอนใยผักหรือหนอนเจาะสมอฝ้าย กัดกินผลึกพิษเข้าไป สภาวะความเป็นด่างสูงในกระเพาะอาหารของแมลง (pH $> 9.0$) จะละลายผลึกโปรตีนและเอนไซม์โปรตีเอสของแมลงจะตัดสายโปรตีนให้กลายเป็นสารพิษที่ออกฤทธิ์แทงทะลุเซลล์เยื่อบุทางเดินอาหาร ทำให้แมลงหยุดกินอาหาร เกิดภาวะโลหิตเป็นพิษและตายภายใน 2--3 วัน
    \item \textbf{เชื้อราก่อโรคแมลง บิวเวอเรีย (\textit{Beauveria bassiana}) และเมตาไรเซียม (\textit{Metarhizium anisopliae})} สปอร์ (Conidia) ของเชื้อราจะเกาะติดกับผนังลำตัว (Cuticle) ของแมลง งอกท่อแทงทะลุเปลือกไคตินเข้าสู่ช่องลำตัว (Hemolymph) เพิ่มจำนวนเส้นใยและปล่อยสารพิษ Beauvericin หรือ Destruxin ทำลายอวัยวะภายใน เมื่อแมลงตาย เส้นใยราสีขาวหรือสีเขียวจะแทงทะลุซากแมลงออกมาภายนอก
\end{enumerate}

\begin{agribox}{เทคนิคการขยายเชื้อสดไตรโคเดอร์มาบนเมล็ดข้าวสุก}
การเตรียมเชื้อราไตรโคเดอร์มาชนิดสดเพื่อใช้ในไร่สวน ให้ใช้ข้าวสารเจ้าผสมน้ำในอัตราส่วน 3 ต่อ 2 หุงด้วยหม้อหุงข้าวไฟฟ้าเมื่อข้าวดีดให้ดงต่อ 10 นาที ตักข้าวขณะยังร้อนใส่ถุงพลาสติกทนร้อนขนาด 6x9 นิ้ว ถุงละ 250 กรัม พับปากถุงทิ้งให้เย็นสนิท จากนั้นหยดหัวเชื้อบริสุทธิ์ 2--3 หยด เขย่าให้เชื้อกระจายทั่ว เจาะรูระบายอากาศใต้หนังยางรัดปากถุง 15--20 รู บ่มที่อุณหภูมิห้องในที่ร่มมีแสงสว่างรำไรเป็นเวลา 7 วัน จะได้สปอร์สีเขียวเข้มพร้อมใช้งาน ผสมน้ำอัตรา 50--100 กรัมต่อน้ำ 20 ลิตร ฉีดพ่นลงดินหรือทรงพุ่ม
\end{agribox}

\section*{คำถามทบทวนท้ายบทที่ 4}
\addcontentsline{toc}{section}{คำถามทบทวนท้ายบทที่ 4}
\begin{enumerate}
    \item จงอธิบายกลไก Mycoparasitism ของเชื้อราไตรโคเดอร์มา และระบุเอนไซม์สำคัญที่ใช้ย่อยผนังเซลล์ของเชื้อราก่อโรค
    \item สารพิษ Cry protein จากแบคทีเรียบีทีมีกลไกการทำลายแมลงศัตรูพืชอย่างไร และเหตุใดจึงมีความปลอดภัยสูงต่อมนุษย์และสัตว์เลือดอุ่น
    \item จงเปรียบเทียบข้อแตกต่างระหว่างการชักนำความต้านทานแบบ ISR (Induced Systemic Resistance) กับ SAR (Systemic Acquired Resistance)
    \item เชื้อ \textit{Streptomyces} sp. SR1301 ที่แยกจากพื้นที่ปกปักทรัพยากร อพ.สธ. มรภ.รำไพพรรณี มีคุณลักษณะและศักยภาพเด่นในการนำมาประยุกต์เป็นชีวภัณฑ์ควบคุมโรคพืชอย่างไร
\end{enumerate}

\clearpage